% 附录 A Design File Management for Synthesis
\bichapter{面向综合的设计文件管理}{Design File Management for Synthesis}
\label{app:filemgmt}

设计（设计描述）存储在设计文件中。设计文件必须具有唯一名称。
若设计是层次化的，每个子设计引用另一个设计文件，该文件也必须具有唯一名称。
注意：不同设计文件可包含同名子设计。

本章包括以下节：
\begin{itemize}
  \item 管理设计数据
  \item 面向综合的划分（Partitioning）
  \item 面向综合的 HDL 编码
\end{itemize}

% ============================================================
\section{管理设计数据}
\subsection*{Managing the Design Data}

使用系统化的组织方法管理设计数据。管理设计数据的两个基本要素是：
\begin{itemize}
  \item 设计数据控制
  \item 设计数据组织
\end{itemize}

% ------------------------------------------------------------
\subsection{设计数据控制}
\subsubsection*{Design Data Control}

随着设计新版本的创建，必须维护某种归档与记录方法，以提供设计演进历史，并在数据丢失时能够重新启动设计过程。
建立数据创建、维护、覆盖与删除的控制是基本的设计管理问题。建立文件命名约定是数据创建最重要的规则之一。
表~\ref{tab:file-ext} 列出各设计数据类型的推荐文件扩展名。

\begin{table}[htbp]
\centering
\caption{文件名扩展名}
\label{tab:file-ext}
\begin{tabularx}{\textwidth}{@{}l l X@{}}
\toprule
\textbf{设计数据类型} & \textbf{扩展名} & \textbf{说明} \\
\midrule
设计源代码 & \texttt{.v} & Verilog \\
 & \texttt{.vhd} & VHDL \\
 & \texttt{.con} & Constraint \\
综合脚本 & \texttt{.scr} & Script \\
报告与日志 & \texttt{.rpt} & Report \\
 & \texttt{.log} & Log \\
设计数据库 & \texttt{.ddc} & Synopsys 内部数据库格式 \\
\bottomrule
\end{tabularx}
\end{table}

% ------------------------------------------------------------
\subsection{设计数据组织}
\subsubsection*{Design Data Organization}

建立并遵守数据组织方法，比选择何种方法更重要。
在将基本设计数据置于一致控制之下后，即可创建有意义的数据组织。为简化数据交换与搜索，设计者应遵守该数据组织系统。

可用层次目录结构解决数据组织问题。您的 compile 策略会影响目录结构。
图~\ref{fig:topdown-dir} 给出基于自顶向下 compile 策略的目录结构，图~\ref{fig:bottomup-dir} 给出自底向上 compile 策略。

\figplaceholder{Figure 82: Top-Down Compile Directory Structure}{自顶向下 Compile 目录结构}{fig:topdown-dir}

\figplaceholder{Figure 83: Bottom-Up Compile Directory Structure}{自底向上 Compile 目录结构}{fig:bottomup-dir}

\begin{seeAlsoBox}
Compile 策略
\end{seeAlsoBox}

% ============================================================
\section{面向综合的划分}
\subsection*{Partitioning for Synthesis}

有效划分设计可提升综合结果、缩短 compile 时间，并简化约束与脚本文件。

划分影响块大小；尽管 DC Explorer 没有固有的块大小限制，仍应谨慎控制块大小。
若块过小，可能造成限制有效优化的人为边界；若创建过大块，compile 运行时间可能很长。

使用下列策略划分设计并改善优化与运行时间：
\begin{itemize}
  \item 为设计重用而划分
  \item 将相关组合逻辑放在一起
  \item 对块输出寄存器化
  \item 按设计目标划分
  \item 按 compile 技术划分
  \item 将可共享资源放在一起
  \item 将用户定义资源与其所驱动的逻辑放在一起
  \item 隔离特殊功能，例如 pad、时钟、边界扫描与异步逻辑
\end{itemize}

下列各节描述每种策略。

% ------------------------------------------------------------
\subsection{为设计重用而划分}
\subsubsection*{Partitioning for Design Reuse}

设计重用通过减少设计、集成与测试工作来缩短上市时间。

重用现有设计时，应划分设计以便实例化这些设计。

为使设计可被重用，在划分与块设计期间遵循下列指南：
\begin{itemize}
  \item 彻底定义并文档化设计接口
  \item 尽可能标准化接口
  \item 对 HDL 代码参数化
\end{itemize}

% ------------------------------------------------------------
\subsection{将相关组合逻辑放在一起}
\subsubsection*{Keeping Related Combinational Logic Together}

默认情况下，DC Explorer 不能将逻辑移过层次边界。将相关组合逻辑分到不同块会引入限制逻辑优化的人为障碍。

为获得最佳结果，应用下列策略：
\begin{itemize}
  \item 将相关组合逻辑及其目标寄存器分组在一起。\\
        在处理完整组合路径时，DC Explorer 有灵活性合并逻辑，从而得到更小、更快的设计。
        将组合逻辑与其 destination register 分组还可简化 timing constraint，并支持 sequential optimization。
  \item 消除粘合逻辑（glue logic）。\\
        粘合逻辑是连接各块的组合逻辑。将此逻辑移入某一块可通过为 DC Explorer 提供额外灵活性来改善综合结果。
        消除粘合逻辑也缩短 compile 时间，因为 DC Explorer 需要优化的逻辑级数更少。
\end{itemize}

例如，假设设计在关键路径上或附近包含三个组合云。图~\ref{fig:poor-part} 给出该设计的不良划分。
每个组合云出现在单独块中，因此 DC Explorer 无法充分发挥其组合优化技术。

\figplaceholder{Figure 84: Poor Partitioning of Related Logic}{相关逻辑的不良划分}{fig:poor-part}

图~\ref{fig:good-part} 给出同一设计且无人为边界的情况。在此设计中，DC Explorer 有灵活性在组合云中合并相关功能。

\figplaceholder{Figure 85: Keeping Related Logic in the Same Block}{将相关逻辑保持在同一块中}{fig:good-part}

% ------------------------------------------------------------
\subsection{对块输出寄存器化}
\subsubsection*{Registering Block Outputs}

为简化约束定义，确保寄存器驱动块输出，如图~\ref{fig:reg-outputs} 所示。

\figplaceholder{Figure 86: Registering All Outputs}{对所有输出寄存器化}{fig:reg-outputs}

\begin{lstlisting}
set_driving_cell my_flop [all_inputs]
set_input_delay 2 -clock CLK
\end{lstlisting}

该方法使您能轻松约束每个块，因为：
\begin{itemize}
  \item 单个块输入上的驱动强度始终等于平均输入驱动的驱动强度
  \item 来自前一块的输入延时始终等于触发器的路径延时
\end{itemize}

由于所有输出均已寄存器化时不存在仅组合路径，对设计做时间预算并使用 \cmd{set\_output\_delay} 命令更容易。
鉴于每个模块内发生一个时钟周期，各模块的约束简单且相同。

此划分方法可改善仿真性能。所有输出均已寄存器化时，可用仅边沿触发的进程描述模块。
敏感列表仅包含时钟，或许还有复位引脚。有限的敏感列表通过使进程在每个时钟周期仅触发一次来加速仿真。

% ------------------------------------------------------------
\subsection{按设计目标划分}
\subsubsection*{Partitioning by Design Goal}

将具有不同设计目标的逻辑划分到不同块。当设计某些部分比其他部分更面积与时序关键时使用此方法。

为获得最佳综合结果，将非关键速度约束逻辑与关键速度约束逻辑隔离。
通过隔离非关键逻辑，可对块施加不同约束，例如最大面积约束。

图~\ref{fig:diff-goals} 给出如何分离具有不同设计目标的逻辑。

\figplaceholder{Figure 87: Blocks With Different Constraints}{具有不同约束的块}{fig:diff-goals}

% ------------------------------------------------------------
\subsection{按 Compile 技术划分}
\subsubsection*{Partitioning by Compile Technique}

将需要不同 compile 技术的逻辑划分到不同块。当设计同时包含高度结构化逻辑与随机逻辑时使用此方法。
\begin{itemize}
  \item highly structured logic（如通常包含大型 XOR tree 的 error-detection circuitry）更适合 structuring。
  \item 随机逻辑更适合 flattening。
\end{itemize}

% ------------------------------------------------------------
\subsection{将可共享资源放在一起}
\subsubsection*{Keeping Sharable Resources Together}

DC Explorer 可共享大型 resource（如 adder 或 multiplier），但仅当这些 resource 属于同一 VHDL \texttt{process} 或 Verilog \texttt{always} block 时，才能进行 resource sharing。

例如，若两个独立加法器具有相同目标路径且对其路径有多路复用输出，请将加法器保持在一个 VHDL process 或 Verilog always 块中。
若约束允许共享，此方法允许 DC Explorer 共享资源（使用一个加法器而非两个）。图~\ref{fig:share-res} 给出逻辑示例的可能实现。

\figplaceholder{Figure 88: Keeping Sharable Resources in the Same Process}{将可共享资源保持在同一进程中}{fig:share-res}

关于资源共享的更多信息，见 HDL Compiler 文档。

% ------------------------------------------------------------
\subsection{将用户定义资源与其所驱动的逻辑放在一起}
\subsubsection*{Keeping User-Defined Resources With the Logic They Drive}

用户定义资源是用户定义的函数、过程或宏单元，或用户创建的 DesignWare 组件。
DC Explorer 不能自动共享或创建用户定义资源的多个实例。
但将这些资源与其所驱动的逻辑放在一起，可在单次实例化无法达到时序目标时，通过手动插入用户定义资源的多次实例化来拆分负载。

图~\ref{fig:dup-udr} 说明当信号 \texttt{PARITY\_ERR} 上的负载过大以致无法满足约束时，通过多次实例化拆分负载。

\figplaceholder{Figure 89: Duplicating User-Defined Resources}{复制用户定义资源}{fig:dup-udr}

% ------------------------------------------------------------
\subsection{隔离特殊功能}
\subsubsection*{Isolating Special Functions}

将特殊功能（如 I/O pad、时钟生成电路、边界扫描逻辑与异步逻辑）与核心逻辑隔离。
图~\ref{fig:top-part} 给出设计顶层的推荐划分。

\figplaceholder{Figure 90: Recommended Top-Level Partitioning}{推荐的顶层划分}{fig:top-part}

设计顶层包含 I/O pad 环与中间层次；中间层次包含边界扫描逻辑、时钟生成电路、异步逻辑与核心逻辑的子模块。
中间层次的存在是为了允许实例化 I/O pad 的灵活性。隔离时钟生成电路可实现该模块的实例化与仔细仿真。
隔离异步逻辑有助于将可测性问题与静态时序分析问题限制在小范围内。

% ============================================================
\section{面向综合的 HDL 编码}
\subsection*{HDL Coding for Synthesis}

HDL 编码是综合的基础，因为它隐含设计的初始结构。编写 HDL 源代码时，始终考虑代码的硬件含义。
良好的编码风格可生成更小、更快的设计。本节提供信息以帮助编写高效代码，从而在最短时间内达到设计目标。

主题包括：
\begin{itemize}
  \item 编写与工艺无关的 HDL
  \item 使用 HDL 构造
  \item 编写有效代码
\end{itemize}

% ------------------------------------------------------------
\subsection{编写与工艺无关的 HDL}
\subsubsection*{Writing Technology-Independent HDL}

使用完全自动综合流程的高层次设计目标是：不实例化任何门或触发器。
若达到此目标，您将拥有可读、简洁且可移植的高层次 HDL 代码，可转移到其他厂商或未来工艺。

在某些情况下，HDL Compiler 工具需要编译器指令以提供实现信息，同时仍保持与工艺无关。
在 Verilog 中，编译器指令以字符 \texttt{//} 或 \texttt{/*} 开头。在 VHDL 中，编译器指令以两个连字符（\texttt{--}）后跟 \texttt{pragma} 或 \texttt{synopsys} 开头。更多信息见 HDL Compiler 文档。

下列各节讨论保持 HDL 代码与工艺无关的各种方法。

\subsubsection{推断组件}
\paragraph*{Inferring Components}

HDL Compiler 提供推断下列组件的能力：
\begin{itemize}
  \item 多路选择器（multiplexer）
  \item 寄存器
  \item 三态驱动器
  \item 多位组件（multibit components）
\end{itemize}

这些推断能力在下列各页讨论。更多信息与示例见 HDL Compiler 文档。

\paragraph{推断多路选择器}
\textit{Inferring Multiplexers}

HDL Compiler 可从 HDL 代码中的 \texttt{case} 语句推断通用多路选择器单元（\texttt{MUX\_OP}）。
若目标工艺库至少包含 2 选 1 多路选择器单元，DC Explorer 将推断的 \texttt{MUX\_OP} 映射到目标工艺库中的多路选择器单元。
DC Explorer 在 compile 期间根据设计约束确定 \texttt{MUX\_OP} 实现。

使用 \texttt{infer\_mux} 编译器指令控制多路选择器推断。附着到块时，该指令强制对块中所有 \texttt{case} 语句做多路选择器推断；附着到 \texttt{case} 语句时，强制对该特定 \texttt{case} 语句做多路选择器推断。

\paragraph{推断寄存器}
\textit{Inferring Registers}

寄存器推断允许您在设计中指定与工艺无关的时序逻辑。寄存器是简单的 1 位存储器件，可以是锁存器或触发器。
锁存器是电平敏感存储器件；触发器是边沿触发存储器件。

每当您未在所有条件下指定输出的结果值时（如不完全指定的 \texttt{if} 或 \texttt{case} 语句），HDL Compiler 推断 D 锁存器。
HDL Compiler 也可推断 SR 锁存器与主从锁存器。

每当 Verilog \texttt{always} 块或 VHDL \texttt{process} 的敏感列表包含边沿表达式（测试信号上升或下降沿）时，HDL Compiler 推断 D 触发器。
HDL Compiler 也可推断 JK 触发器与翻转触发器。

\paragraph{混合寄存器类型}
\textit{Mixing Register Types}

为获得最佳结果，将每个 Verilog \texttt{always} 块或 VHDL \texttt{process} 限制为单一类型的寄存器推断：锁存器、带异步置位/复位的锁存器、触发器、带异步置位/复位的触发器，或带同步置位/复位的触发器。

在设计中混合上升沿与下降沿触发的触发器时要小心。
若模块同时推断上升沿与下降沿触发的触发器，且目标工艺库不含下降沿触发的触发器，DC Explorer 会在下降沿时钟的时钟树中生成反相器。

\paragraph{推断无控制信号的寄存器}
\textit{Inferring Registers Without Control Signals}

对于推断无控制信号的寄存器，使数据与时钟引脚可从输入端口或通过组合逻辑控制。
若门级仿真器不能从输入端口或通过组合逻辑控制数据或时钟引脚，仿真器无法初始化电路，仿真将失败。

\paragraph{推断带控制信号的寄存器}
\textit{Inferring Registers With Control Signals}

可用异步或同步控制信号初始化或控制触发器状态。

对于推断锁存器上的异步控制信号，使用 \texttt{async\_set\_reset} 编译器指令（VHDL 中为属性）标识异步控制信号。
推断触发器时，HDL Compiler 自动标识异步控制信号。

对于推断同步复位，使用 \texttt{sync\_set\_reset} 编译器指令（VHDL 中为属性）标识同步控制。

\paragraph{推断三态驱动器}
\textit{Inferring Three-State Drivers}

将高阻值（Verilog 中为 \texttt{1'bz}，VHDL 中为 \texttt{'Z'}）赋给输出引脚，以使 DC Explorer 推断三态门。
三态逻辑降低设计可测性并使调试困难。在可能的情况下，用多路选择器替换三态缓冲器。

切勿在条件表达式中使用高阻值。HDL Compiler 始终将与高阻值比较的表达式求值为假，这可能导致门级实现与 RTL 描述行为不同。

关于三态推断的更多信息，见 HDL Compiler 文档。

\paragraph{推断多位组件}
\textit{Inferring Multibit Components}

多位推断允许您将多路选择器、寄存器与三态驱动器映射到规则结构化逻辑或多位库单元。使用多位组件可带来下列结果：
\begin{itemize}
  \item 因共享晶体管与优化的晶体管级版图而更小的面积与延时
  \item sequential gate 中的 clock skew 减小
  \item sequential banked component 的 clock power 降低
  \item 改善 datapath 的规则版图
\end{itemize}

在下列情况下多位组件可能效率不高：
\begin{itemize}
  \item 作为状态机寄存器
  \item 在更适合 single-bit design 的小型 bused logic 中
\end{itemize}

在决定映射到多位还是单位组件时，必须权衡多位组件的收益与优化灵活性损失。

将 \texttt{infer\_multibit} 编译器指令附着到总线信号以推断多位组件。
也可在优化后用 \cmd{create\_multibit} 与 \cmd{remove\_multibit} 命令在单位与多位实现之间切换。

\subsubsection{设计状态机}
\paragraph*{Designing State Machines}

可用多种不同格式指定状态机：
\begin{itemize}
  \item Verilog
  \item VHDL
  \item 状态表（state table）
  \item PLA
\end{itemize}

若在 HDL 代码中使用 \texttt{state\_vector} 与 \texttt{enum} 编译器指令，DC Explorer 可从网表提取状态表。
在状态表格式中，DC Explorer 不保留 \texttt{casex}、\texttt{casez} 与 \texttt{parallel\_case} 信息。
DC Explorer 不优化无效输入组合与互斥输入。

图~\ref{fig:fsm-arch} 给出有限状态机架构。

\figplaceholder{Figure 91: Finite State Machine Architecture}{有限状态机架构}{fig:fsm-arch}

使用提取的状态表具有下列收益：
\begin{itemize}
  \item 可执行状态最小化
  \item 可在不同编码风格之间权衡
  \item 可在不 flattening 设计的情况下使用 don't care 条件
  \item 自动推导 don't care 状态码
\end{itemize}

% ------------------------------------------------------------
\subsection{使用 HDL 构造}
\subsubsection*{Using HDL Constructs}

下列各节提供关于下列特定 HDL 构造的信息与指南：
\begin{itemize}
  \item 一般 HDL 构造
  \item Verilog 宏定义
  \item VHDL 端口定义
\end{itemize}

\subsubsection{一般 HDL 构造}
\paragraph*{General HDL Constructs}

本节信息同时适用于 Verilog 与 VHDL。

\paragraph{敏感列表}
\textit{Sensitivity Lists}

应为每个 Verilog \texttt{always} 块或 VHDL \texttt{process} 完整指定敏感列表。不完整的敏感列表（如下例所示）可导致 HDL 与门级设计之间的仿真不匹配。

\begin{lstlisting}[caption={不完整的敏感列表（Verilog）},label=ex:sens-v]
always @ (A)
   C <= A | B;
\end{lstlisting}

\begin{lstlisting}[caption={不完整的敏感列表（VHDL）},label=ex:sens-vhdl]
process (A)
   C <= A or B;
\end{lstlisting}

\paragraph{值赋值}
\textit{Value Assignments}

Verilog 与 VHDL 都支持在 RTL 代码中使用立即与延迟值赋值。
立即值赋值——由 Verilog 阻塞赋值（\texttt{=}）与 VHDL 变量（\texttt{:=}）实现——所生成的硬件依赖于赋值顺序。
延迟值赋值——由 Verilog 非阻塞赋值（\texttt{<=}）与 VHDL 信号（\texttt{<=}）实现——所生成的硬件与赋值顺序无关。

为获得正确仿真结果：
\begin{itemize}
  \item 在 sequential Verilog \texttt{always} block 或 VHDL \texttt{process} 内使用 delayed（nonblocking）assignment
  \item 在 combinational Verilog \texttt{always} block 或 VHDL \texttt{process} 内使用 immediate（blocking）assignment
\end{itemize}

\paragraph{if 语句}
\textit{if Statements}

当作为连续赋值一部分用在 Verilog \texttt{always} 块或 VHDL \texttt{process} 中的 \texttt{if} 语句不含 \texttt{else} 子句时，DC Explorer 创建锁存器。
下列示例给出在综合期间生成锁存器的 \texttt{if} 语句。

\begin{lstlisting}[caption={不正确的 if 语句（Verilog）},label=ex:if-v]
if ((a == 1) && (b == 1))
   z = 1;
\end{lstlisting}

\begin{lstlisting}[caption={不正确的 if 语句（VHDL）},label=ex:if-vhdl]
if (a = '1' and b = '1') then
   z <= '1';
end if;
\end{lstlisting}

\paragraph{case 语句}
\textit{case Statements}

若 \texttt{if} 语句包含三个以上条件，考虑使用 \texttt{case} 语句以改善设计并行性与代码清晰度。
下列示例使用 \texttt{case} 语句实现 3 位译码器。

\begin{lstlisting}[caption={使用 case 语句（Verilog）},label=ex:case-v]
case ({a, b, c})
   3'b000: z = 8'b00000001;
   3'b001: z = 8'b00000010;
   3'b010: z = 8'b00000100;
   3'b011: z = 8'b00001000;
   3'b100: z = 8'b00010000;
   3'b101: z = 8'b00100000;
   3'b110: z = 8'b01000000;
   3'b111: z = 8'b10000000;
   default: z = 8'b00000000;
endcase
\end{lstlisting}

\begin{lstlisting}[caption={使用 case 语句（VHDL）},label=ex:case-vhdl]
case_value :=  a & b & c;
CASE case_value IS
   WHEN "000" =>
      z <= "00000001";
   WHEN "001" =>
      z <= "00000010";
   WHEN "010" =>
      z <= "00000100";
   WHEN "011" =>
      z <= "00001000";
   WHEN "100" =>
      z <= "00010000";
   WHEN "101" =>
      z <= "00100000";
   WHEN "110" =>
      z <= "01000000";
   WHEN "111" =>
      z <= "10000000";
   WHEN OTHERS =>
      z <= "00000000";
END CASE;
\end{lstlisting}

不完整的 \texttt{case} 语句会导致创建锁存器。VHDL 不支持不完整的 \texttt{case} 语句。
在 Verilog 中，可通过 \texttt{default} 子句或 \texttt{full\_case} 编译器指令避免锁存器推断。

尽管 \texttt{full\_case} 指令与 \texttt{default} 子句都防止锁存器推断，它们含义不同。
\texttt{full\_case} 指令断言已指定所有有效输入值且无需 \texttt{default} 子句；\texttt{default} 子句为任何未定义输入值指定输出。

为获得最佳结果，使用 \texttt{default} 子句而非 \texttt{full\_case} 指令。
若未指定的输入值是 don't care 条件，使用输出值为 \texttt{x} 的 \texttt{default} 子句可生成更小的实现。

若使用 \texttt{full\_case} 指令，每当 \texttt{case} 表达式求值为未指定输入值时，门级仿真可能与 RTL 仿真不匹配。
若使用 \texttt{default} 子句，仅当指定了 don't care 条件且 \texttt{case} 表达式求值为未指定输入值时，才可能发生仿真不匹配。

\paragraph{常量定义}
\textit{Constant Definitions}

使用 Verilog 的 \texttt{`define} 语句或 VHDL 的 \texttt{constant} 语句定义全局常量。将全局常量定义保存在单独文件中。
使用参数（Verilog）或 generic（VHDL）定义局部常量。

示例~\ref{ex:macro-v} 给出包含全局 \texttt{`define} 与局部参数的 Verilog 代码片段。
示例~\ref{ex:const-vhdl} 给出包含全局常量与局部 generic 的 VHDL 代码片段。

\begin{lstlisting}[caption={使用宏与参数（Verilog）},label=ex:macro-v]
// Define global constant in def_macro.v
`define WIDTH 128

// Use global constant in reg128.v
reg regfile[WIDTH-1:0];

// Define and use local constant in module myfile
module myfile (a, b, c);
   parameter WIDTH=128;
   input [WIDTH-1:0] a, b;
   output [WIDTH-1:0] c;
\end{lstlisting}

\begin{lstlisting}[caption={使用全局常量与 Generics（VHDL）},label=ex:const-vhdl]
-- Define global constant in synthesis_def.vhd
constant WIDTH : INTEGER := 128;

-- Include global constants
library my_lib;
USE my_lib.synthesis_def.all;

-- Use global constant in entity my_design
entity my_design is
   port (a,b : in std_logic_vector(WIDTH-1 downto 0);
         c: out std_logic_vector(WIDTH-1 downto 0));
end my_design;

-- Define and use local constant in entity my_design
entity my_design is
   generic (WIDTH_VAR : INTEGER := 128);
   port (a,b : in std_logic_vector(WIDTH-1 downto 0);
         c: out std_logic_vector(WIDTH-1 downto 0));
end my_design;
\end{lstlisting}

\subsubsection{使用 Verilog 宏定义}
\paragraph*{Using Verilog Macro Definitions}

在 Verilog 中，宏用 \texttt{`define} 语句实现。对 \texttt{`define} 语句遵循下列指南：
\begin{itemize}
  \item 仅用 \texttt{`define} 语句声明常量。
  \item 将 \texttt{`define} 语句保存在单独文件中。
  \item 不要使用嵌套的 \texttt{`define} 语句。\\
        读取嵌套超过两层的宏很困难。为使代码可读，不要使用嵌套的 \texttt{`define} 语句。
  \item 不要在模块定义内使用 \texttt{`define}。\\
        在模块定义内使用 \texttt{`define} 时，局部宏与全局宏具有相同引用名但不同值。使用参数定义局部常量。
\end{itemize}

\subsubsection{使用 VHDL 端口定义}
\paragraph*{Using VHDL Port Definitions}

在 VHDL 源代码中定义端口时，遵循下列指南：
\begin{itemize}
  \item 使用 \texttt{STD\_LOGIC} 与 \texttt{STD\_LOGIC\_VECTOR} 包。\\
        使用 \texttt{STD\_LOGIC} 可避免在综合后设计上需要类型转换函数。
  \item 不要使用 \texttt{buffer} 端口模式。\\
        将 port 声明为 \texttt{buffer} 时，该 port 必须在整个 hierarchy 中均以 \texttt{buffer} 模式使用。
        为简化综合，应将 port 声明为 \texttt{out}，再定义一个驱动该 output port 的内部 signal。
\end{itemize}

% ------------------------------------------------------------
\subsection{编写有效代码}
\subsubsection*{Writing Effective Code}

本节提供为综合编写高效、可读 HDL 源代码的指南。指南涵盖：
\begin{itemize}
  \item 标识符
  \item 表达式
  \item 函数
  \item 模块
\end{itemize}

\subsubsection{标识符指南}
\paragraph*{Guidelines for Identifiers}

好的标识符名称传达信号含义、变量值或模块功能；没有这些信息，硬件描述难以阅读。

遵循下列命名指南以改善 HDL 源代码可读性：
\begin{itemize}
  \item 确保信号名称传达信号含义或变量值，同时不过于冗长。\\
        例如，假设有一个表示 \texttt{rs1} 浮点操作码的变量。短名称如 \texttt{frs1} 不能向读者传达含义；
        长名称如 \texttt{floating\_pt\_opcode\_rs1} 传达含义，但其长度可能使源代码难读。使用如 \texttt{fpop\_rs1} 的名称可同时满足两个目标。
  \item 对大小写以及区分名称中各词使用一致的命名风格。\\
        常用风格包括 C、Pascal 与 Modula。
  \begin{itemize}
    \item C 风格使用小写名称并用下划线分隔词，例如 \texttt{packet\_addr}、\texttt{data\_in}、\texttt{first\_grant\_enable}。
    \item Pascal 风格将名称首字母与每个词首字母大写，例如 \texttt{PacketAddr}、\texttt{DataIn}、\texttt{FirstGrantEnable}。
    \item Modula 风格对名称首字母用小写，并将后续词首字母大写，例如 \texttt{packetAddr}、\texttt{dataIn}、\texttt{firstGrantEnable}。
  \end{itemize}
        选择一种约定并一致应用。
  \item 避免易混淆字符。\\
        某些字符（字母与数字）看起来相似且易混淆，例如 O 与 0（零）；l 与 1（一）。
  \item 避免保留字。
  \item 名称采用 noun 或 noun 后接 verb 的形式，例如 \texttt{AddrDecode}、\texttt{DataGrant}、\texttt{PCI\_interrupt}。
  \item 添加后缀以澄清名称含义。\\
        表~\ref{tab:suffix} 给出常用后缀及其含义。
\end{itemize}

\begin{table}[htbp]
\centering
\caption{信号名称后缀及其含义}
\label{tab:suffix}
\begin{tabularx}{\textwidth}{@{}l X@{}}
\toprule
\textbf{后缀} & \textbf{含义} \\
\midrule
\texttt{\_clk} & 时钟信号 \\
\texttt{\_next} & 寄存之前的信号 \\
\texttt{\_n} & 低有效信号 \\
\texttt{\_z} & 连接到三态输出的信号 \\
\texttt{\_f} & 使用有效下降沿的寄存器 \\
\texttt{\_xi} & 芯片主输入 \\
\texttt{\_xo} & 芯片主输出 \\
\texttt{\_xod} & 芯片主开漏输出 \\
\texttt{\_xz} & 芯片主三态输出 \\
\texttt{\_xbio} & 芯片主双向 I/O \\
\bottomrule
\end{tabularx}
\end{table}

\subsubsection{表达式指南}
\paragraph*{Guidelines for Expressions}

对表达式遵循下列指南：
\begin{itemize}
  \item 使用括号表示优先级。\\
        表达式运算符优先级规则令人困惑，因此应使用括号使表达式易于阅读。
        除非使用 DesignWare 资源，括号对生成逻辑影响很小。难以阅读且无括号的逻辑表达式示例：
\begin{lstlisting}
bus_select = a ^ b & c~^d|b^~e&^f[1:0];
\end{lstlisting}
  \item 用函数调用或连续赋值替换重复表达式。\\
        若特定表达式使用超过两三次，考虑用实现该表达式的函数或连续赋值替换该表达式。
\end{itemize}

\subsubsection{函数指南}
\paragraph*{Guidelines for Functions}

对函数遵循下列指南：
\begin{itemize}
  \item 不要在函数内使用全局引用。\\
        在过程代码中，函数在被调用时求值。在连续赋值中，当其任一已声明输入改变时求值函数。\\
        避免在函数内引用全局名，因为全局值改变时函数可能不会被重新求值。这可能导致 HDL 描述与门级网表之间的仿真不匹配。\\
        例如，下列 Verilog 函数引用全局名 \texttt{byte\_sel}：
\begin{lstlisting}
function byte_compare;
   input [15:0] vector1, vector2;
   input [7:0] length;

   begin
      if (byte_sel)
         // compare the upper byte
      else
         // compare the lower byte
      ...
   end
endfunction // byte_compare
\end{lstlisting}
  \item 注意任务与函数的局部存储是静态的。\\
        形式参数、输出与局部变量在函数返回后保留其值。每次调用函数时重用局部存储。
        此存储对调试有用，但存储重用也意味着函数与任务不能递归调用。
  \item 使用 component implication 时要谨慎。\\
        可用 \texttt{map\_to\_module} 与 \texttt{return\_port\_name} 编译器指令将函数映射到特定实现。
        simulation 使用 function 内容；synthesis 则以 gate-level module 替代该 function。使用 component implication 时，RTL model 与 gate-level model 可能不同。
        因此，在对门级设计运行仿真之前，无法完全验证设计。\\
        下列功能可能需要 component instantiation 或 functional implication：
  \begin{itemize}
    \item 用于省电的时钟门控电路
    \item 具有潜在冒险的异步逻辑（包括异步逻辑以及在某些状态下有效的异步信号）
    \item Datapath 电路（包括大型多路选择器；实例化的宽多路选择器组；存储器元件如 RAM 或 ROM；以及黑盒宏单元）
  \end{itemize}
        关于 component implication 的更多信息，见 HDL Compiler 文档。
\end{itemize}

\subsubsection{模块指南}
\paragraph*{Guidelines for Modules}

对模块遵循下列指南：
\begin{itemize}
  \item 通过端口传递值时避免使用逻辑表达式。\\
        端口列表可包含表达式，但表达式使调试复杂化。此外，隔离与位域相关的问题很困难，尤其当该位域导致内部端口量与外部端口量不同时。
  \item 将局部引用定义为 generic（VHDL）或参数（Verilog）。不要将 generic 或参数传入模块。
\end{itemize}
